\section{Conclusiones}

El periodo de practicas en ActLab, dentro del proyecto \textit{Didascalia VC}, ha permitido consolidar una participacion tecnica completa: desde diagnostico y optimizacion inicial de rendimiento VR, hasta dise\~no de sistemas extensibles de comportamiento, refuerzo de la capa de conexion con web, ampliacion funcional de conflictos y validacion final en Meta Quest 3.

La principal aportacion ha sido contribuir a que la aplicacion gane en tres dimensiones clave:
\begin{itemize}[leftmargin=1.5em]
    \item \textbf{Fiabilidad}: mejor control de errores, aserciones y trazabilidad ante eventos inesperados.
    \item \textbf{Escalabilidad}: pipeline de conflictos y estructura de comportamiento preparada para crecimiento.
    \item \textbf{Operatividad}: despliegue funcional en dispositivo objetivo y soporte a pruebas en multiples servidores.
\end{itemize}

Desde el punto de vista formativo, la experiencia ha servido para aplicar de forma integrada conocimientos del grado en un entorno profesional real, con restricciones de producto, necesidades de equipo y exigencia de entrega continua.

Como cierre administrativo para la entrega en GIPE, solo resta completar los datos identificativos y la firma de tutorizacion empresarial, sin que ello afecte al contenido tecnico ya documentado.

\vspace{1cm}
\noindent\textbf{Firma del tutor en la empresa:}\\[1.5cm]
\noindent Nombre y cargo: \textbf{[COMPLETAR]}\\
\noindent Firma y sello: \textbf{[COMPLETAR]}\\
\noindent Fecha: \textbf{[COMPLETAR]}